% =============================================================================
%  latex-agentic — article template
%  Engine: XeLaTeX ONLY (compile with `latexmk` or `make pdf`).
%
%  Why XeLaTeX: native UTF-8 input and system/OpenType fonts via fontspec.
%  Do NOT add \usepackage{inputenc} or \usepackage{fontenc} — they are
%  pdfLaTeX-era packages and will fight fontspec. Just type UTF-8 directly.
% =============================================================================
\documentclass[11pt,a4paper]{article}

% --- Fonts & encoding (XeLaTeX) --------------------------------------------
% fontspec must load before most font-related packages. We keep the default
% Latin Modern fonts (shipped with TeX Live) so this template compiles
% anywhere; swap in any installed font with \setmainfont if you like.
\usepackage{fontspec}
\defaultfontfeatures{Ligatures=TeX}

% --- Language ---------------------------------------------------------------
% polyglossia is the XeLaTeX-native replacement for babel.
\usepackage{polyglossia}
\setdefaultlanguage{english}

% --- Typography -------------------------------------------------------------
\usepackage{microtype}        % subtle protrusion/expansion — better grey balance
\usepackage[autostyle=true]{csquotes}  % context-aware quotation marks
\usepackage{geometry}
\geometry{margin=2.5cm}
\usepackage{setspace}
\setstretch{1.06}

% --- Tables, figures, math --------------------------------------------------
\usepackage{booktabs}         % proper rules: \toprule \midrule \bottomrule
\usepackage{graphicx}
\usepackage{mwe}              % ships example-image* placeholder graphics
\usepackage{amsmath,amssymb}
\usepackage{caption}
\captionsetup{font=small,labelfont=bf}
\usepackage{xcolor}           % needed for the blue!55!black link colour below

% --- Bibliography (biblatex + biber) ----------------------------------------
% authoryear style; bibliographic database lives in refs.bib.
\usepackage[backend=biber,style=authoryear,sorting=nyt,maxbibnames=10]{biblatex}
\addbibresource{refs.bib}

% --- Cross-references -------------------------------------------------------
% hyperref must be loaded LAST (after almost everything), and cleveref must be
% loaded AFTER hyperref. This ordering avoids the classic "hyperref redefines
% \ref" breakage. cleveref gives \cref / \Cref which print "Section 2",
% "Figure 1", etc. automatically.
\usepackage[
  unicode=true,
  colorlinks=true,
  linkcolor=black,
  citecolor=black,
  urlcolor=blue!55!black,
  pdftitle={A Concise Study of Reproducible Typesetting},
  pdfauthor={Ada A. Researcher},
  pdfsubject={Sample XeLaTeX article from latex-agentic},
  pdfkeywords={LaTeX, XeLaTeX, reproducibility, typesetting}
]{hyperref}
\usepackage{cleveref}

% --- Metadata ---------------------------------------------------------------
% The new-article skill replaces these three sample values (the title, the
% author line, and the date) with the resolved title/author/date. The sample
% strings below are chosen to be unique so an exact-match edit is unambiguous;
% if you are editing by hand, just overwrite them.
\title{A Concise Study of Reproducible Typesetting}
\author{Ada A. Researcher\\\small Department of Document Engineering, Example University}
\date{\today}

\begin{document}
\maketitle

\begin{abstract}
\noindent This template demonstrates a complete, genuinely good-looking
XeLaTeX article: a structured body, a \texttt{booktabs} table, a figure, a
numbered equation referenced with \texttt{cleveref}, and an
\texttt{authoryear} bibliography driven by \texttt{biber}. The sample text
below is filler intended to be replaced; the surrounding scaffolding is meant
to be kept.
\end{abstract}

% ---------------------------------------------------------------------------
\section{Introduction}\label{sec:intro}
Reproducible typesetting matters whenever a document must be rebuilt months
later from source. Early discussions of the topic framed it as a question of
toolchain stability~\autocite{lamport1994latex}, while later work emphasised
the role of consistent fonts and encodings~\autocite{robertson2023fontspec}.
This article walks through the pieces a working author actually needs:
sectioning, floats, mathematics, and citations. \Cref{sec:results} presents a
small table and figure; \cref{sec:math} states the one equation we rely on.

% ---------------------------------------------------------------------------
\section{Results}\label{sec:results}
\Cref{tab:summary} reports illustrative measurements, and \cref{fig:overview}
sketches the corresponding layout. Both are placeholders: replace the numbers
and the graphic with your own.

\begin{table}[htbp]
  \centering
  \caption{Illustrative build timings across three engines (sample data).}%
  \label{tab:summary}
  \begin{tabular}{lrr}
    \toprule
    Engine    & Cold build (s) & Warm build (s) \\
    \midrule
    pdfLaTeX  & 4.1            & 1.2            \\
    XeLaTeX   & 5.3            & 1.6            \\
    LuaLaTeX  & 6.0            & 1.8            \\
    \bottomrule
  \end{tabular}
\end{table}

\begin{figure}[htbp]
  \centering
  \includegraphics[width=0.6\linewidth]{example-image}
  \caption{Placeholder graphic supplied by the \texttt{mwe} package. Swap in
  your own figure with \texttt{\textbackslash{}includegraphics}.}%
  \label{fig:overview}
\end{figure}

% ---------------------------------------------------------------------------
\section{A Touch of Mathematics}\label{sec:math}
For completeness we include one numbered equation. The Gaussian integral,
shown in \cref{eq:gauss}, is a convenient sanity check that math mode and
cross-referencing both work:
\begin{equation}
  \label{eq:gauss}
  \int_{-\infty}^{\infty} e^{-x^{2}}\,\mathrm{d}x = \sqrt{\pi}.
\end{equation}
The reader is invited to confirm \cref{eq:gauss} numerically.

% ---------------------------------------------------------------------------
\section{Conclusion}\label{sec:conclusion}
We have assembled the minimum viable article: structure, a float pair, a
referenced equation, and a working bibliography. The same skeleton scales to a
full paper~\autocite{knuth1986texbook}; replace the prose, keep the
preamble. The micro-typographic refinements enabled by \texttt{microtype}
have a measurable perceptual effect~\autocite{nguyen2022microtype}. For a
survey of online tooling, see~\autocite{overleaf2024guide}.

% ---------------------------------------------------------------------------
\printbibliography[heading=bibintoc,title={References}]

\end{document}
